\documentclass{article}

\usepackage{amsmath} % math stuff
\usepackage{amssymb} % math stuff
\usepackage{array} % equations and stuff
\usepackage{bm} % bold math
%\usepackage{booktabs} % extra table rule options
%\usepackage{caption} % suppressed table numbering; incompatible with revtex, and longtable, I think
\usepackage{comment} % comment environment
%\usepackage{enumitem} % customization of enumeration, itemize, and description
\usepackage[T1]{fontenc} % font encoding for special characters, must also use scalable font package
\usepackage[margin=0.8in]{geometry} % paper sizes and margins (but be careful not to mess up pre-defined pages)
\usepackage{graphicx} % for graphics
%\usepackage{helvet} % default font is the helvetica postscript font
\usepackage[utf8]{inputenc} % special characters in tex input
\usepackage{layouts} % print units like widths
\usepackage{lipsum} % lorem ipsum filler text
\usepackage{lmodern} % scalable font?
\usepackage{longtable} % multi-page tables
\usepackage{makecell} % specify line-breaks in table cells
\usepackage{mathrsfs} % math script font
\usepackage{mhchem} % easier chemical formula
\usepackage{microtype} % allows disabling of ligatures
\usepackage{multicol} % multicolumns
%\usepackage{newcent} % new century schoolbook font
\usepackage{nicefrac}
\usepackage{numprint} % print and format (large) numbers
\usepackage{parskip} % removes paragraph indentation, and adjusts paragraph skip, as well as list items
\usepackage{pdfpages} % add pdf files as pages
%\usepackage{setspace} % adjust text spacing and indents
\usepackage{siunitx} % decimal alignment, and degree symbol
\usepackage{subfigure} % divided figures
%\usepackage{tabu} % extra table options
\usepackage{textcomp} % symbols
\usepackage{threeparttablex} % better footnotes with longtable
\usepackage{titling} % title placement
\usepackage{ulem} % strikethrough text
\usepackage{upgreek} % upright Greek letters
%\usepackage{url} % superceded by hyperref
\usepackage{verbatim} % verbatim environment
\usepackage{xcolor} % colors and color boxes
\usepackage{xspace} % commands that don't eat up white space
\usepackage{hyperref} % links and page setup; should always come last

\hypersetup{
 bookmarks=true,
 colorlinks=true,
 citecolor=blue,
 linkcolor=blue,
 urlcolor=blue,
 pdfstartview={XYZ null null 1.0} % default open view is 100%
}

\DisableLigatures[f,t]{encoding = T1} % disable ff, fi, fl, tt ligatures; without options, it also disables -- = endash
\renewcommand{\arraystretch}{1.0} % extra vertical (and horizontal?) space in tables

% define centered, left- and right-aligned columns with specified widths
\newcommand{\PreserveBackslash}[1]{\let\temp=\\#1\let\\=\temp}
\newcolumntype{C}[1]{>{\PreserveBackslash\centering}p{#1}}
\newcolumntype{L}[1]{>{\PreserveBackslash\raggedright}p{#1}}
\newcolumntype{R}[1]{>{\PreserveBackslash\raggedleft}p{#1}}

\begin{document}

\pagestyle{empty} % don't number pages

% custom title
\begin{center}
{\LARGE Express Riddler}

\vspace{0.15in}

{\Large 6 May 2022}
\end{center}


\section*{Riddle:}

You are on the 10th floor of a tower and want to exit on the first floor.
You get into the elevator and hit 1.
However, this elevator is malfunctioning in a specific way.
When you hit 1, it correctly registers the request to descend, but it randomly selects some floor below your current floor (including the first floor).
The car then stops at that floor.
If it's not the first floor, you again hit 1 and the process repeats.

Assuming you are the only passenger on the elevator, how many floors on average will it stop at (including your final stop, the first floor) until you exit?


\section*{Solution:}

I let $E_{n}$ be the expected number of stops to get from the $(n+1)^{\textrm{th}}$ floor to the first floor.
So, getting from the tenth floor to the first floor would require $E_{9}$ stops on average.

Trivially, $E_{0}=0$ and $E_{1}=1$, because starting from the second floor, there is only the possibility of stopping on the first floor after a single stop.
Starting from the $(n+1)^{\textrm{th}}$ floor, it takes a single stop to get to some new floor $k+1$, and from that first stop, the expected number of additional stops is $E_{k}$ for $k<n$.
For each floor below $n+1$, there is a \nicefrac{1}{n} chance of the first stop being that floor.
So $E_{n}$ can be written recursively as

\begin{align*}
E_{n} &= 1+\frac{1}{n}(E_{0}+\dots+E_{n-1}) \\
        &= 1+\frac{1}{n}(E_{1}+\dots+E_{n-1}) \\
        &= 1+\frac{1}{n}\sum_{k=1}^{n-1}E_{k}
\end{align*}

Rearraging and adjusting the index of this equation gives

\[
\sum_{k=1}^{n}E_{k} = (n+1)\left(E_{n+1}-\frac{1}{n+1}\right)-n
\]

The harmonic numbers are defined by

\[
H_{n+1} = H_{n}+\frac{1}{n+1}
\]

with $H_{1}=1$.
Another known identity of the harmonics numbers is

\begin{align*}
\sum_{k=1}^{n}H_{k} &= (n+1)H_{n}-n \\
                    &= (n+1)\left(H_{n+1}-\frac{1}{n+1}\right)-n
\end{align*}

Thus, $H_{n}$ and $E_{n}$ are defined identically, so $E_{n}=H_{n}$.
The ninth harmonic number $H_{9}$ is $\nicefrac{7129}{2520}\approx2.829$, so the solution is \fcolorbox{red}{white}{\textbf{\nicefrac{7129}{2520}}}\,.


\end{document}